\chapter{Ingesta de datos}
\label{sec:man-ingesta}

\section{Formatos soportados}

\begin{itemize}
    \item \textbf{Sysmon}: EVTX exportado o XML de
    \code{Microsoft-Windows-Sysmon/Operational}. IDs 1, 3, 7, 8,
    10, 11, 13, 22, 23.
    \item \textbf{Sentinel / Log Analytics}: JSON de queries KQL o
    stream vía conector. Tablas: \code{SigninLogs},
    \code{AuditLogs}, \code{SecurityEvent},
    \code{DeviceProcessEvents}, \code{AzureActivity},
    \code{SecurityAlert}.
    \item \textbf{FortiAnalyzer}: CSV con esquema propietario
    (\code{date}, \code{time}, \code{srcip}, \code{dstip}, etc.).
    \item \textbf{AWS Cognito}: JSON de CloudWatch Logs para
    eventos de Cognito.
    \item \textbf{IIS}: W3C Extended Log con \code{\#Fields}.
    \item \textbf{CSV genérico}: cualquier CSV con
    \code{FieldConfig} mapeado.
    \item \textbf{Generic (texto)}: fallback heurístico para
    formatos sin parser específico.
\end{itemize}

Para detalle del mapeo de cada formato a entidades del grafo, ver
capítulo \emph{Formatos de log} en el libro de Especificación.

\section{One-shot vs streaming}

\textbf{One-shot}: subís un archivo (drag-and-drop, o
\emph{Ingesta \textrightarrow{} Upload}). El sistema lo procesa
íntegro, construye el grafo, cierra. Óptimo para investigaciones
sobre logs históricos.

\textbf{Streaming}: el conector de Sentinel abre una query KQL
continua; los eventos llegan en batches y se agregan al grafo
vivo. Óptimo para monitoreo activo o construcción progresiva de
una ventana de análisis.

Elegir en \emph{Ingesta \textrightarrow{} Modo}.

\section{Field mapping}

Todo formato termina produciendo un \code{FieldConfig} — un objeto
que mapea campos del log a roles semánticos del grafo:

\begin{itemize}
    \item \code{TIME} — timestamp (obligatorio).
    \item \code{ACTOR} — quién hizo la acción (usuario, proceso).
    \item \code{ACTION} — qué acción fue (autenticar, conectar).
    \item \code{TARGET} — sobre qué recurrió (host, archivo, IP).
    \item \code{INFO} — atributos adicionales (puerto, bytes).
\end{itemize}

Para formatos conocidos (Sysmon, Sentinel, Forti, Cognito, IIS),
el \code{FieldConfig} viene preconfigurado. Para CSV genérico, la
UI ofrece preview y asignación de roles:

\begin{enumerate}
    \item Arrastrás el archivo.
    \item Ves un preview de las primeras 100 filas con los roles
    propuestos.
    \item Podés ajustar cada columna a otro rol con el dropdown.
    \item Confirmar \textrightarrow{} el sistema aplica el mapping
    e ingesta.
\end{enumerate}

\section{Ingestor asistido por LLM}

Para formatos no reconocidos (y para evolución de formatos
conocidos), GraphHunter incluye el carril \emph{agéntico}. El
flujo típico:

\begin{enumerate}
    \item Si el sistema no reconoce el formato,
    \emph{Ingesta \textrightarrow{} Asistente \textrightarrow{}
    Generar mapping}.
    \item El LLM lee una muestra y propone un
    \code{FieldConfig}.
    \item El preview muestra la propuesta; revisás los roles y
    mapeo a entidades.
    \item Aprobar \textrightarrow{} el mapping se compila a VRL y
    queda disponible.
    \item Próxima ingesta del mismo formato lo usa
    automáticamente.
\end{enumerate}

El LLM nunca ingesta eventos en caliente. Propone datos, el
compilador determinístico emite código, el runtime lo ejecuta. Ver
libro de Arquitectura (cap.\ \emph{Ingestor agéntico}) y
Especificación (cap.\ \emph{Formatos de log}) para el detalle.

\section{Archivos grandes}

Para archivos de decenas de GB (EVTX de dominio grande,
exportación de varios meses de Sentinel):

\begin{itemize}
    \item \textbf{Preferir streaming} al modo one-shot.
    \item Activar \emph{Ingesta \textrightarrow{} Opciones
    \textrightarrow{} Ingesta por disco} si el archivo supera el RAM
    disponible. El sistema usa un store intermedio en disco.
    \item Ajustar \code{GRAPHHUNTER_INGEST_BATCH_SIZE} a 5000-20000
    según memoria.
    \item Para casos extremos (100 GB+), considerar partir el
    archivo antes con herramientas externas y ingestar por partes a
    la misma sesión.
\end{itemize}

\section{Dead Letter Queue}

Eventos que se parsean pero violan invariantes (timestamp nulo,
entidad faltante, relación inválida) van a la DLQ — no se pierden,
no corrompen el grafo. Para revisar:

\begin{itemize}
    \item \emph{Panel \textrightarrow{} DLQ} muestra todos los
    eventos pendientes con la razón de rechazo.
    \item Corregido el mapping (o el evento), \emph{Reingerir}
    vuelve a intentar.
    \item \emph{Purgar} vacía la cola sin reintentar.
\end{itemize}

Una DLQ llena es señal de mapping mal configurado. Antes de
purgar, inspeccionar qué eventos fallan y por qué.

\section{Drift}

Cuando un formato que ingerís hace tiempo cambia aguas arriba (un
campo nuevo, un separador distinto), el monitor de drift lo
detecta y abre una propuesta de \emph{re-review} en el panel
agéntico. El sistema sigue ingestando con el mapping actual
mientras tanto. Revisar y aprobar el mapping actualizado cuando te
venga bien.
